\scanpage{39}\mspage{38}
\[
  (S',T')\preccurlyeq(S,T)
\]
signifie que $S'=S$, et $T'$ plus fine que $T$ i.e. l'application identique $S\xrightarrow{\mathrm{id}}S$ est continue.
\margin{Description de $\bigwedge_i(S_i,T_i)$, c'est $S=\bigcap S_i$ muni de la \inserted{\uuline{ultra}}\uline{topologie} \deleted{ultra} \uline{initiale} sur $S$ pour \inserted{la famille des} $S\hookrightarrow S_i$ (muni de $T_i$ : \deleted{cette ultra}top. pour une famille quelc. d'applications \deleted{\unsure{$f$}} $f_i\colon S\mathbin{-}S_i$}

Revenons à la situation (on reprend les notations $F\ldots$)
\diagram{\begin{tikzcd}[ampersand replacement=\&,column sep=small,row sep=tiny]
G \& {\subset} \& F\\
\rotatebox{90}{$\preccurlyeq$}\\
\deleted{\unsure{F}}G'
\end{tikzcd}}
donc, posant $T=|G|$, $S=|F|$
\margin{de $S$ dans des $S_i$ munis de $T_i$ \unsure{existe} toujours, \unsure{en|on} \uline{ultr.} \uline{topologie} \uline{engendrée} par les \inserted{fermés} $f_i^{-1}(F_i)$ ($F_i$ fermé ds $T_i$) : on prend les \deleted{\unsure{réunions}} \inserted{intersections} quelc. de réunions quelc. d'ensembles de la forme $f_i^{-1}(F_i)\ldots$}
\diagram{\begin{tikzcd}[ampersand replacement=\&,column sep=small,row sep=tiny]
\deleted{\unsure{\text{on a}|\text{elle}}} \& T\arrow[r,hook,no head] \& S\\
\& \tau_0^{+} \& \tau\\
\& \rotatebox{90}{$\leq$}\\
\& \tau'_0
\end{tikzcd}}
$T$ fermé \unsure{de|dans} $(S,\tau)$, $\tau_0$ top. induite par $\tau$, $\tau'_0$ top. plus fine. On considère les \deleted{topologies} $F'\preccurlyeq F$ qui induisent $G'$, i.e. les topologies $\tau'$ sur $S$ \unsure{telles|telle} que $\tau'\leq\tau$ (plus fines que $\tau$ -- il y a \uline{plus} de fermés) et qui \unsure{\deleted{in}induisent|induisent} $\tau'_0$.
\margin{correspond à la notion de relation de préordre \deleted{\unsure{engendrée}} \inserted{induite} \unsure{par} \deleted{\unsure{des}} \inserted{famille} \inserted{images inverses} \unsure{(}relations de préordres) : \deleted{\unsure{donc ici}} $x\leq y\Longleftrightarrow f_i(x)\leq f_i(y)$ \ $\forall i\in I$}
Donc on veut que
\[
\begin{aligned}
 1^\circ\quad &(S,\tau')\longrightarrow(S,\tau)\quad\text{continue},\\
 \text{et}\qquad 2^\circ)\quad &(\deleted{\unsure{S}}T,\tau'_0)\longrightarrow(S,\tau')\quad\text{continue}
\end{aligned}
\]
situation mixte à cheval sur celle des top. initiales et finales \ldots
\diagram{\begin{tikzcd}[ampersand replacement=\&,column sep=large,row sep=tiny]
T\arrow[r,"{i}"] \& S\arrow[r,"\mathrm{id}"] \& S\\
\tau'_0 \& \tau' ? \& \tau
\end{tikzcd}}
\margin{on sait que le composé est continu.}

Donc on veut (i) que les $A\in\unsure{P|\mathcal P|\mathscr P}(S)$ fermés pour $\tau$, \unsure{i.e.} $A\in P_\tau(S)$, soient fermés pour $\tau'$, \deleted{i.e.} et (ii) que pour \uline{$\forall$} $A'$ fermé pour $\tau'$, $i^{-1}(A')=A'\cap T$ soit fermé pour $\tau'_0$. \inserted{Si on veut \unsure{$=$'|$\Leftarrow$} $\tau'_0=\tau'\mid T$,} $T$ étant fermé \deleted{\unsure{pou}} pour $\tau$ donc pour $\tau'$, les parties \inserted{\unsure{fermées|fermés}} \unsure{de|\deleted{de}} $T$ \inserted{pour $\tau'_0$} sont \unsure{fermées|fermés}
